\subsection{Differential characterization of etale algebras}

THEOREM 3. --- Let A be a commutative algebra of finite degree over K. For A to be etale it is necessary and sufficient that the module \(\Omega_{\mathrm{K}}(\mathrm{A})\) of K-differentials of A should be reduced to 0.

\begin{enumerate}
\def\labelenumi{\Alph{enumi})}
\item
  Let L be an algebraic closure of K (V, p.~23, Th. 2). For A to be etale it is necessary and sufficient that the algebra \(\mathbf{A}_{(L)}\) over L should be diagonalizable (V, p.~30, Prop. 2). Further, the A-module \(\Omega_{\mathrm{L}}(\mathbf{A}_{(L)})\) is isomorphic to \(\Omega_{\mathrm{K}}(\mathbf{A}) \otimes_{\mathbf{A}} \mathbf{A}_{(L)}\) (III, p.~572, Prop. 20), hence to \(\Omega_{\mathrm{K}}(\mathbf{A}) \otimes_{\mathbf{K}} \mathbf{L}\) , by the associativity of the tensor product; therefore \(\Omega_{\mathrm{K}}(\mathbf{A}) = 0\) is equivalent to \(\Omega_{\mathrm{L}}(\mathbf{A}_{(L)}) = 0\) . To prove Th. 3 it is therefore enough to consider the case where K is algebraically closed and to show that the algebra A is diagonalizable if and only if \(\Omega_{\mathrm{K}}(\mathbf{A}) = 0\) .
\item
  Suppose that A is diagonalizable; then (V, p.~29, Prop. 1), the vector space A is generated by the idempotents of A. The assertion \(\Omega_{\mathrm{K}}(\mathbf{A}) = 0\) is thus a consequence of the following lemma:
\end{enumerate}

Lemma 2. --- Let A be a commutative algebra over K and e an idempotent of A; then we have de = 0 in \(\Omega_{\mathbf{K}}(\mathbf{A})\) .

From the relation \(e = e^2\) we deduce \(de = 2e \cdot de\) ; on multiplying by \(e\) we obtain \(e \cdot de = 2e \cdot de\) , hence \(e \cdot de = 0\) , and so finally, \(de = 2e \cdot de = 0\) .

\begin{enumerate}
\def\labelenumi{\Alph{enumi})}
\setcounter{enumi}{2}
\tightlist
\item
  We first prove two lemmas :
\end{enumerate}

Lemma 3. --- Let A be a commutative algebra of finite degree over the algebraically closed field K, such that \(\Omega_{\mathbf{K}}(\mathbf{A}) = 0\) . Then we have \(\mathfrak{m} = \mathfrak{m}^2\) for every maximal ideal \(\mathfrak{m}\) of A.

The algebra \(\mathbf{A} / \mathfrak{m}\) is an extension of finite degree of the algebraically closed field \(\mathbf{K}\) , whence \([\mathbf{A} / \mathfrak{m}:\mathbf{K}] = 1\) . Hence for each \(a\in \mathbf{A}\) there exists a unique scalar \(\lambda\) such that \(a - \lambda ,1\in \mathfrak{m}\) ; write \(\mathrm{D}(a)\) for the residue class of \(a - \lambda .1\) mod \(\mathfrak{m}^2\) . It is clear that \(\mathrm{D}\) is a K-derivation of \(\mathbf{A}\) into the A-module \(\mathfrak{m} / \mathfrak{m}^2\) . The universal property of \(\Omega_{\mathbf{K}}(\mathbf{A})\) (III, p.~569) and the hypothesis \(\Omega_{\mathbf{K}}(\mathbf{A}) = 0\) now imply that \(\mathrm{D} = 0\) , whence \(\mathfrak{m} / \mathfrak{m}^2 = 0\) and so \(\mathfrak{m} = \mathfrak{m}^2\) .

Lemma 4. --- Let A be a commutative ring and let \(\mathfrak{a}\) be a finitely generated ideal of A such that \(\mathfrak{a} = \mathfrak{a}^2\) . Then there exists an idempotent \(e\) in A such that \(\mathfrak{a} = \mathbf{A}e\) .

Let \((a_{1},\ldots ,a_{r})\) be a generating system of the ideal \(\mathfrak{a}\) ; since \(\mathfrak{a} = \mathfrak{a}^2\) , there exist elements \(x_{ij}\) in \(\mathfrak{a}\) such that \(a_{i} = \sum_{j = 1}^{r}x_{ij}a_{j}\) for \(1\leqslant i\leqslant r\) . Write M for the square

matrix of order r whose elements are \(\delta_{ij} - x_{ij}\) and let D be its determinant. There exists (III, p.~532, Formula (26)) a square matrix N of order r with elements in A such that N . M = D . I \(_{r}\) , whence immediately Da \(_{j}\) = 0 for \(1 \leqslant j \leqslant r\) and so finally Da = 0. Now the matrix M is congruent to I \(_{r}\) mod a, hence D ≡ 1 mod a. Put e = 1 - D ; then \(e \in a\) and ex = x for all \(x \in a\) . It follows that e is an idempotent and a is equal to Ae.

With these lemmas established, let us show by induction on the degree of A that A is diagonalizable if K is algebraically closed and \(\Omega_{\mathrm{K}}(\mathrm{A})=0\) . Let m be a maximal ideal of A (I, p.~104). By Lemmas 3 and 4, there exists an idempotent e such that m = Ae; we have seen that A/m is of degree 1 over K. Hence A is a direct sum of the ideals \(\mathfrak{a}=(1-e)\mathbf{A}\) and m and we have \([\mathfrak{a}:K]=1\) , hence A is isomorphic to \(K\times A/a\) . Since \(\Omega_{\mathrm{K}}(\mathrm{A}/\mathfrak{a})\) is isomorphic to a quotient of \(\Omega_{\mathrm{K}}(\mathrm{A})\) (III, p.~573, Prop. 22), it is zero and the induction hypothesis shows that A/ \(\mathfrak{a}\) is diagonalizable. This then shows A to be diagonalizable.

